\subsection{Every adjacent norm in the complex seventeen tower}

\begin{statement}{E.64}{Proved}{Explicit global norm images, completed Hilbert 90, and the next fixed class group}
Keep the actual first cocycle $A=3\kappa_3(17)$ of E.62.
Put $K=\mathbb Q(\zeta)$, $\zeta=\zeta_3$, and choose
\[
 L_d=K(\beta_d),\qquad
 \beta_0=17\zeta,\qquad \beta_{d+1}^3=\beta_d
 \quad(d\geq0).
 \tag{E64.1}
\]
Thus $L_0=K$ and $\beta_d^{3^d}=17\zeta$.
Write $\widehat{L_d^\times}_3=\varprojlim_m
 L_d^\times/(L_d^\times)^{3^m}$, with its prodiscrete topology.
The following formulas evaluate the image and every nonempty fiber of
each actual adjacent norm
\[
 N_d:\widehat{L_{d+1}^\times}_3
                   \longrightarrow\widehat{L_d^\times}_3.
 \tag{E64.2}
\]

There is exactly one prime $\mathfrak t_d$ over $17$ and exactly one
prime $\mathfrak l_d$ over $3$ in $L_d$. Their residue fields are
$\mathbb F_{289}$ and $\mathbb F_3$, respectively. The extension
$L_{d+1}/L_d$ is cyclic of degree three, totally ramified at these two
primes and unramified elsewhere. Its specified generator is
$\sigma_d(\beta_{d+1})=\zeta\beta_{d+1}$.
For a prime $\mathfrak q\nmid 3\cdot17$ of $L_d$, define
$\gamma_{d,\mathfrak q}\in\mathbb F_3$ by the actual finite-field equation
\[
 \overline{\beta_d}^{\,(\mathrm N\mathfrak q-1)/3}
       =\overline\zeta^{\,\gamma_{d,\mathfrak q}},
 \qquad
 \mathcal I_d=\{\mathfrak q\nmid3\cdot17:
                         \gamma_{d,\mathfrak q}\ne0\}.
 \tag{E64.3}
\]
Precisely the primes in $\mathcal I_d$ are inert in
$L_{d+1}/L_d$; every other prime away from $3,17$ splits completely.

For $z\in L_d^\times$, let $v_{\mathfrak t_d}(\beta_d)=1$ and set
\[
 \begin{split}
 \delta_{d,17}(z)&=
  \log_{\bar\zeta}
  \left(\overline{z\beta_d^{-v_{\mathfrak t_d}(z)}}^{\,96}\right)
                         \ \in\mathbb F_3,\\
 \delta_{d,\mathfrak q}(z)&=
       \gamma_{d,\mathfrak q}v_{\mathfrak q}(z)\pmod3
                   \quad(\mathfrak q\in\mathcal I_d),\\
 \Theta_d(z)&=(\delta_{d,17}(z),
                 (\delta_{d,\mathfrak q}(z))_{\mathfrak q\in\mathcal I_d}).
 \end{split}
 \tag{E64.4}
\]
The bar in the first line is reduction of the displayed unit at
$\mathfrak t_d$. Every coordinate kills cubes, so the same map on the
completion is evaluated on any representative of its class modulo
cubes. The sequence of coordinates away from $17$ has finite support.
There is an exact sequence of topological abelian groups
\[
 \widehat{L_{d+1}^\times}_3
 \xrightarrow{\,N_d\,}\widehat{L_d^\times}_3
 \xrightarrow{\,\Theta_d\,}
 \mathbb F_3\ \oplus\
                  \bigoplus_{\mathfrak q\in\mathcal I_d}\mathbb F_3
 \longrightarrow0 .
 \tag{E64.5}
\]
The right-hand group has the discrete topology. In particular the
complete global norm test uses only the explicit residue computations
and valuation congruences in (E64.4); its image is an open subgroup.
The omitted local norm condition at the single prime over $3$ follows
from global reciprocity, as proved below.

Both exact completed comparisons are
\[
 \begin{split}
 \ker N_d&=(\sigma_d-1)\widehat{L_{d+1}^\times}_3,\\
 (\widehat{L_{d+1}^\times}_3)^{\sigma_d}
        &=\widehat{L_d^\times}_3
               \quad\hbox{under the actual field inclusion},\\
 \widehat{L_{d+1}^\times}_3/
                \widehat{L_d^\times}_3
       &\xrightarrow{\ \sim\ }\ker N_d,\qquad
          u\longmapsto\sigma_d(u)/u .
 \end{split}
 \tag{E64.6}
\]
These are equalities and a topological isomorphism, rather than a
replacement of ordinary multiplicative groups by their completions
without an exactness proof.
For completeness, each nonempty fiber of (E64.2) has the explicit
noncanonical topological form
\[
 \mu_3\times
 \left\{(a_i)_{i\ge1}\in\prod_{i\ge1}\mathbb Z_3:
       \text{for every }m,\ a_i\in3^m\mathbb Z_3
                         \text{ for all but finitely many }i\right\}.
 \tag{E64.7}
\]
Choose a lift $y_0$ of the target and a free basis $u_i$ of the ordinary
norm-one group modulo its roots of unity; the map is
$ (\zeta^e,(a_i))\mapsto
       y_0\zeta^e\prod_i u_i^{a_i}$.
The proof gives the precise basis choices and explains convergence.

At the first degree-eighteen extension there is also the exact
class-group calculation
\[
 \left|\mathrm {Cl}(L_2)^{\mathrm {Gal}(L_2/L_1)}\right|=27,
 \qquad
 E_{L_1}\cap N_{L_2/L_1}L_2^\times=E_{L_1},
 \tag{E64.8}
\]
where $E_{L_i}$ denotes the ordinary unit group.
The actual comparison with unit norms is the exact sequence
\[
 1\longrightarrow \mathrm {Am}_{\rm st}(L_2/L_1)
 \longrightarrow \mathrm {Cl}(L_2)^{\mathrm {Gal}(L_2/L_1)}
 \longrightarrow E_{L_1}/N_{L_2/L_1}E_{L_2}
 \longrightarrow1 .
 \tag{E64.9}
\]
Thus (E64.8) proves element-norm surjectivity on base units.
Formula (E64.9) retains the additional quotient needed to determine
which of those norms have an ordinary unit as preimage.

Under the exact Kummer--Shapiro map E.46, (E64.2) is the complex-factor
component of the actual positional transition
\[
 \begin{array}{ccc}
 \widehat{F_{d+1}^\times}_3\times
       \widehat{L_{d+1}^\times}_3
 &\xrightarrow{\ (N,N_d)\ }&
 \widehat{F_d^\times}_3\times\widehat{L_d^\times}_3,\\[-1mm]
 F_d=\mathbb Q(17^{1/3^d})&&
 \end{array}
 \tag{E64.10}
\]
from level $d+2$ to level $d+1$ of E.62.
Total norms to $\mathbb Q$ commute with this arrow; consequently,
once both component preimages have been chosen, the target's actual
$J$-condition modulo cubes also holds for the source.
The next global question is the existence and evaluation of compatible
choices through all these arrows, together with the pure-factor norm
conditions. An assertion that membership in each individual norm image
already supplies compatible choices is not used.
The manuscript anchor remains \S51.16, scans172--178, folios733v--736v,
equations(371)--(379), and Part E(c),(d).
\end{statement}

\begin{proof}
\emph{Fields and the two ramified primes.}
The polynomial $X^{3^d}-17\zeta$ is Eisenstein at the prime of $K$
over $17$. It proves $[L_d:K]=3^d$, uniqueness of $\mathfrak t_d$,
and the assertions about its uniformizer and residue field.
For $d\ge1$ the absolute minimal polynomial is
\[
 P_d(X)=X^{2\cdot3^d}+17X^{3^d}+289.
 \tag{E64.11}
\]
Modulo three it is $(X+1)^{2\cdot3^d}$, whereas
$P_d(-1)=273=3\cdot91$. Hence $P_d(Y-1)$ is Eisenstein at three.
It proves that $L_d$ has exactly one prime over three, with absolute
ramification index $2\cdot3^d$ and residue degree one.
The statement for $d=0$ is the elementary decomposition in
$\mathbb Q(\zeta)$.
These indices show that both primes are ramified of degree three
in each adjacent extension.
Since $\zeta\in L_d$, adjoining $\beta_{d+1}$ gives the specified
cyclic extension and automorphism.
The relative polynomial $X^3-\beta_d$ has discriminant
$-27\beta_d^2$; away from the two displayed primes it defines an
etale algebra. There is no infinite ramification because $L_d$ is
totally imaginary. This proves the ramification assertions.

\emph{The exact local tests and their global sufficiency.}
At $\mathfrak q\nmid3\cdot17$, the extension is unramified.
The polynomial $X^3-\bar\beta_d$ either splits into three distinct
linear factors or is irreducible, since the residue field contains
$\mu_3$. Its splitting is exactly the zero condition in (E64.3).
For an unramified cubic extension the unit norm is onto and the
norm of a uniformizer has valuation three.
One may prove the unit assertion by first taking the surjective
finite-field norm and then the surjective trace on each successive
principal-unit quotient. Thus the local norm condition is precisely
$\delta_{d,\mathfrak q}=0$ at inert primes.

At $\mathfrak t_d$ the extension is tamely and totally ramified of
degree three. Its norm on residue units is cubing.
Cubing is an automorphism of the principal-unit pro-seventeen group,
and the norm of an element of the base field is its cube.
These facts prove sufficiency as well as necessity of the
residue-unit test. The literal identity
\[
                 N_d(\beta_{d+1})=\beta_d
 \tag{E64.12}
\]
allows every valuation, and leaves exactly the first line of
(E64.4).

Now let $z\in L_d^\times$ have all displayed coordinates zero.
It is a norm at every place other than possibly $\mathfrak l_d$.
Global reciprocity makes the product of its local Artin symbols one.
Since this is the only remaining place, its Artin symbol there is
also one. The local reciprocity kernel is the local norm group.
The Hasse norm theorem for a cyclic extension then gives
$z=N_d(y)$ for an actual $y\in L_{d+1}^\times$.
Here the exact theorem is \cite[VIII, Theorem3.1]{E_MilneCFT},
including the proof by the injection
$L_d^\times/N_dL_{d+1}^\times
\to\bigoplus_vL_{d,v}^\times/N_dL_{d+1,v}^\times$.
Necessity follows from the local calculations.

We also need surjectivity of the entire last map of (E64.5),
rather than just the zero criterion.
For a cyclic extension the relative Brauer exact sequence in
\cite[VIII, \S4, the exact sequence preceding Lemma4.1 and its proof;
IV, Theorem3.14]{E_MilneCFT} identifies the global norm quotient
with the tuples of local norm classes having total invariant zero.
All nonzero local norm quotients here have order three.
Arbitrary finitely supported prescribed classes outside $\mathfrak l_d$
can be supplemented by the unique class at $\mathfrak l_d$ whose
invariant makes their sum zero. This is permitted because its local
degree is three. The relative exact sequence supplies a global class,
and the cyclic crossed-product description represents it by an
element of $L_d^\times$. Each coordinate in (E64.4) is an explicit
isomorphism from the corresponding local norm quotient to
$\mathbb F_3$; a possible change of sign of a local invariant is
an invertible change of that one coordinate.
This proves the asserted surjectivity with exactly the displayed
coordinates.

\emph{Exactness of completion, with both kernels checked.}
We give the required abelian-group argument in additive notation.
If
\[
                 0\longrightarrow C\longrightarrow B
                      \longrightarrow D\longrightarrow0
 \tag{E64.13}
\]
is exact and the three-primary torsion of $D$ is killed by $3^r$,
then
\[
 C\cap3^{n+r}B\subseteq3^nC\subseteq C\cap3^nB .
 \tag{E64.14}
\]
Indeed $c=3^{n+r}b\in C$ implies that the image of $b$ in $D$
is killed by $3^{n+r}$, hence by $3^r$; therefore $3^rb\in C$.
The kernels $C/(C\cap3^nB)$ in the quotient exact sequences have
surjective transition maps. Taking inverse limits is consequently
exact, and (E64.14) identifies the induced topology with completion
for multiplication by $3^n$. This proves exactness of the completed
version of (E64.13), including its injective left-hand map.

Apply this first to
\[
 0\longrightarrow K_d\longrightarrow L_{d+1}^\times
                  \xrightarrow{N_d} I_d^{\rm norm}\longrightarrow0,
 \quad
 K_d=\ker(N_d:L_{d+1}^\times\to L_d^\times),\quad
 I_d^{\rm norm}=N_dL_{d+1}^\times.
 \tag{E64.15}
\]
The torsion of $I_d^{\rm norm}$ is contained in the finite roots of
unity of $L_d$, so (E64.14) applies.
To compare its completion with the target completion, note separately
that $3L_d^\times\subset I_d^{\rm norm}$ in additive notation,
because the norm of a base element is its cube. Consequently
\[
 I_d^{\rm norm}\cap3^{n+1}L_d^\times
             \subseteq3^nI_d^{\rm norm},\qquad
 3^nI_d^{\rm norm}\subseteq
             I_d^{\rm norm}\cap3^nL_d^\times .
 \tag{E64.16}
\]
Its induced and intrinsic filtrations are cofinal.
The exact quotient sequences then identify
$\widehat I_d^{\rm norm}$ with the kernel of the completed
map to $L_d^\times/I_d^{\rm norm}$.
This quotient has exponent three and is exactly the discrete group
evaluated above. It proves (E64.5) on completions and identifies
the completed norm kernel with $\widehat K_d$.

For the second direction use ordinary Hilbert90, in its exact form
\cite[II, Proposition1.22 and Corollary1.23]{E_MilneCFT}:
\[
 0\longrightarrow L_d^\times\longrightarrow L_{d+1}^\times
                  \xrightarrow{\ \sigma_d-1\ }K_d\longrightarrow0 .
 \tag{E64.17}
\]
The torsion of $K_d$ is contained in the finite roots of unity of
$L_{d+1}$, so (E64.14) applies to this different exact sequence too.
Its completed sequence identifies the kernel of $\sigma_d-1$
with the actual included $\widehat{L_d^\times}_3$, and its image
with $\widehat K_d$. The latter has already been identified with
the norm kernel by (E64.15)--(E64.16).
This proves every assertion in (E64.6).
The quotient maps are open: on the intrinsic power filtrations their
images of $3^n$ times the completed source are $3^n$ times the
completed image. Together with the cofinal inclusions this proves
the claimed topological statements.

There is also an explicit way to produce a completed norm preimage
once the test has passed. Represent the target $a$ modulo cubes by
$a_0\in L_d^\times$. The established ordinary criterion supplies
$y_0$ with $N_dy_0=a_0$. For example, an enumeration of the elements
of the explicit number field, followed by exact norm computation,
finds such a $y_0$ and terminates by the theorem.
Write $a/a_0=u^3$ in $\widehat{L_d^\times}_3$. This division exists:
using units and a basis of the principal-divisor group identifies
the multiplicative completion with a finite root-of-unity group
and restricted three-adic sequences; its kernel modulo cubes is
exactly its subgroup of cubes.
Then $y_0u$, with $u$ included from the base field, is the required
preimage. This construction takes no inverse limit over
noncompact nonempty fibers.

\emph{The fibers as explicit topological groups.}
The three-primary roots of unity in every $L_d$ are exactly $\mu_3$.
Indeed, for any cyclotomic three-power extension $C/K$, the polynomial
$X^{3^d}-17\zeta$ remains Eisenstein at a prime over $17$, because
$C/K$ is unramified there. Its degree remains $3^d$, which proves
that $L_d$ and $C$ are linearly disjoint over $K$.
In particular $\zeta_9\notin L_d$.
Since $N_d(\zeta)=1$, the three-primary torsion of $K_d$ is exactly
$\mu_3$.

The quotient of a number field's multiplicative group by its roots
of unity is free abelian of countable rank: use the ordinary unit
group and the subgroup of the free divisor group consisting of
principal divisors. Hence $K_d$ modulo its torsion is free abelian
of at most countable rank. Its rank is infinite.
For infinitely many primes of $L_d$ split in $L_{d+1}$, choose a
prime $\mathfrak P$ above each and an element $v$ with
$(v)=\mathfrak P^h$, where $h$ is the class number of $L_{d+1}$.
The elements $\sigma_d(v)/v$ have nonzero divisors
$h(\sigma_d\mathfrak P-\mathfrak P)$ on disjoint prime orbits;
they are independent. Infinitude of split primes follows from
the proved Chebotarev theorem
\cite[VIII, Theorem7.4]{E_MilneCFT}.
Choose a basis of this free norm-one group modulo its torsion
and lifts $u_i\in K_d$. Since the quotient is free, these choices
split its ordinary torsion subgroup. Hilbert90 gives actual
$v_i$ with $u_i=\sigma_d(v_i)/v_i$.
Completing this direct sum gives exactly (E64.7):
the restrictedness condition ensures that the displayed product
converges modulo every power of three and that its coordinates
are unique. The infinite quotients in this completion are
discrete and need not be finite; the resulting global group
is not being treated as compact.

\emph{The fixed class group in the second complex layer.}
E.57 certified
$\mathrm {Cl}(L_1)\simeq C_3^2$ and the fundamental units
\[
 \begin{split}
 \epsilon_1&=-18+3\beta_1+3\beta_1^2
       -18\beta_1^3/17-5\beta_1^4/17+2\beta_1^5/17,\\
 \epsilon_2&=-18-5\beta_1+2\beta_1^2
       -\beta_1^3/17-8\beta_1^4/17-\beta_1^5/17 .
 \end{split}
 \tag{E64.18}
\]
At the unique prime over $17$, their residues are both
$-1-\bar\zeta=\bar\zeta^2$.
The roots of unity are generated by $-1$ and $\zeta$; each
of these residues is a cube in $\mathbb F_{289}^\times$.
The unit valuations at all unramified primes vanish.
The norm criterion therefore proves that every ordinary unit of
$L_1$ is the norm of an element of $L_2^\times$.
There are exactly two ramified finite primes and no ramified
infinite primes. The ambiguous class-number formula
\cite[Theorem1, formula(1)]{E_LemmermeyerAmbiguous} now gives
\[
 |\mathrm {Cl}(L_2)^{\mathrm {Gal}(L_2/L_1)}|
     =9\,\frac{3^{2-1}}{[E_{L_1}:E_{L_1}\cap N_dL_2^\times]}
     =27 .
 \tag{E64.19}
\]
The same source's Proposition1 gives (E64.9).
Its actual map sends an ambiguous ideal class represented by
$\mathfrak a$, with $\sigma_d\mathfrak a=(\alpha)\mathfrak a$,
to $N_d\alpha$ modulo norms of units. This formula specifies
exactly the additional unit-norm comparison.
No full class-group or fundamental-unit calculation in the
degree-eighteen field is needed for (E64.19).

The retained certificate verifies the exact prime decompositions
and field inclusions at degrees two, six and eighteen, certifies
the degree-six class group and units, and compares (E64.3) with
actual prime contraction at all $58$ primes tested.
Its independent Python part checks $1{,}508$ finite-field norm
identities, $124$ characteristic-zero norm identities and
six Eisenstein expansions. The theorem's statements for every
$d$ follow from the proofs above; finite tests are accompanied
by the exact field and norm maps.
\end{proof}
